% ============================================================
\chapter{Applications}
\label{ch:applications}
% ============================================================

\begin{intuition}
Reinforcement learning has moved from theoretical curiosity to
transformative technology. From superhuman game play to robotic
manipulation, autonomous driving, recommendation systems, and large
language model alignment, RL is at the core of some of the most
impactful AI systems. This chapter surveys key application domains,
highlighting the RL techniques and challenges specific to each.
\end{intuition}

% ------------------------------------------------------------
\section{Game Playing}
\label{sec:app:games}
% ------------------------------------------------------------

\subsection{Board Games: Go, Chess, and Shogi}

\begin{definition}[AlphaGo and AlphaZero]
\index{AlphaGo}\index{AlphaZero}
AlphaGo (Silver et al., 2016) combined supervised learning from expert
games with RL self-play and MCTS. AlphaZero (2017) removed the need for
human data entirely, learning from self-play alone:
\begin{itemize}
    \item \textbf{Representation}: the board is encoded as a multi-channel
    image (stone positions, liberties, history).
    \item \textbf{Network}: a ResNet outputs both a policy $\mathbf{p}$
    and a value $v$ from the board state.
    \item \textbf{Training}: MCTS generates training targets; the network
    is trained to predict MCTS policy and game outcome.
    \item \textbf{Result}: superhuman in Go, chess, and shogi with a
    single algorithm and architecture.
\end{itemize}
\end{definition}

\begin{proposition}[Self-Play Improvement]
In two-player zero-sum games, self-play generates a curriculum of
increasing difficulty. Each iteration of training against the current
best policy produces a stronger opponent, driving continuous improvement
without external data.
\end{proposition}

\subsection{Video Games: Atari and StarCraft}

\begin{example}[Atari (DQN)]
DQN (Mnih et al., 2015) demonstrated that a single deep RL algorithm
can learn directly from pixels across 49 Atari games, achieving
superhuman performance on the majority. Rainbow (2018) pushed
performance further by combining multiple DQN improvements.
\end{example}

\begin{example}[StarCraft II (AlphaStar)]
\index{AlphaStar}
AlphaStar (Vinyals et al., 2019) achieved Grandmaster level in
StarCraft II, a real-time strategy game with:
\begin{itemize}
    \item Imperfect information (fog of war)
    \item Enormous action space ($\sim 10^{26}$ possible actions per step)
    \item Long-horizon planning (games last $\sim$20 minutes)
\end{itemize}
Key techniques: multi-agent league training, Transformer-based
architecture, pointer networks for action selection.
\end{example}

% ------------------------------------------------------------
\section{Robotics}
\label{sec:app:robotics}
% ------------------------------------------------------------

\begin{definition}[Sim-to-Real Transfer]
\index{sim-to-real}
\textbf{Sim-to-real transfer} trains RL policies in simulation and
deploys them on physical robots. The \textbf{reality gap} (mismatch
between simulation and reality) is addressed via:
\begin{itemize}
    \item \textbf{Domain randomization}: randomize physics parameters,
    textures, lighting, and noise during training.
    \item \textbf{System identification}: calibrate the simulator to
    match the real system.
    \item \textbf{Domain adaptation}: fine-tune in the real world with
    few interactions.
\end{itemize}
\end{definition}

\begin{example}[Dexterous Manipulation]
OpenAI demonstrated solving a Rubik's cube with a robotic hand (2019)
using PPO with massive domain randomization (thousands of physics
parameters varied during training). The policy transferred to the
real robot without any real-world training.
\end{example}

\begin{example}[Locomotion]
Legged locomotion has been a major success of deep RL:
\begin{itemize}
    \item Anymal (ETH Zurich): PPO-trained quadruped traversing
    challenging terrain.
    \item Stanford Pupper / Unitree Go1: sim-to-real locomotion policies
    using teacher-student distillation.
    \item Key challenge: robustness to perturbations, terrain variation,
    and hardware damage.
\end{itemize}
\end{example}

\begin{attention}[title=\textbf{Safety in Robotics}]
RL for robotics requires constrained and safe learning (Chapter
\ref{ch:constrained}). Physical hardware is expensive and fragile.
Safety constraints during training (not just at convergence) are
essential to prevent damage.
\end{attention}

% ------------------------------------------------------------
\section{Autonomous Driving}
\label{sec:app:driving}
% ------------------------------------------------------------

\begin{definition}[RL for Autonomous Driving]
\index{autonomous driving}
RL approaches to autonomous driving typically address:
\begin{itemize}
    \item \textbf{Decision making}: lane changes, merging, intersection
    navigation (often modeled as MDPs or POMDPs).
    \item \textbf{Motion planning}: generating smooth, collision-free
    trajectories.
    \item \textbf{End-to-end driving}: mapping sensor inputs directly to
    steering/acceleration commands.
\end{itemize}
\end{definition}

\begin{remark}
Pure RL for autonomous driving faces challenges: safety-critical
requirements, rare but important events (edge cases), and the need for
interpretable decisions. In practice, RL is combined with rule-based
safety layers, imitation learning, and planning.
\end{remark}

\begin{example}[Wayve and Tesla]
Industry applications include:
\begin{itemize}
    \item \textbf{Wayve}: end-to-end RL-based driving in urban
    environments with sim-to-real transfer.
    \item \textbf{Tesla}: uses imitation learning from human drivers
    combined with RL for decision optimization in Autopilot/FSD.
\end{itemize}
\end{example}

% ------------------------------------------------------------
\section{Recommendation Systems}
\label{sec:app:recommender}
% ------------------------------------------------------------

\begin{definition}[RL for Recommendations]
\index{recommendation systems}
A recommendation system can be modeled as an MDP:
\begin{itemize}
    \item \textbf{State}: user interaction history, context, user profile.
    \item \textbf{Action}: item to recommend (from a large catalog).
    \item \textbf{Reward}: click, purchase, engagement time, long-term
    user satisfaction.
\end{itemize}
\end{definition}

\begin{proposition}[Long-Term vs Short-Term Optimization]
Greedy recommendation (maximize immediate click probability) often leads
to filter bubbles and user fatigue. RL enables optimization of
long-term engagement by modeling multi-step interactions and exploring
diverse content.
\end{proposition}

\begin{example}[Industry Applications]
\begin{itemize}
    \item \textbf{YouTube}: uses RL to optimize watch time while
    balancing user satisfaction and content diversity (REINFORCE-based).
    \item \textbf{Netflix}: offline RL methods to optimize
    session-level engagement from logged interaction data.
    \item \textbf{E-commerce}: multi-armed bandits and contextual
    bandits for product ranking and pricing.
\end{itemize}
\end{example}

% ------------------------------------------------------------
\section{LLM Alignment: RLHF and Beyond}
\label{sec:app:llm}
% ------------------------------------------------------------

\begin{definition}[LLM Alignment]
\index{LLM alignment}
\textbf{Alignment} refers to training large language models to be
helpful, harmless, and honest (the 3H criteria). RLHF (Chapter
\ref{ch:constrained}) is the dominant approach:
\begin{enumerate}
    \item Collect human preference data on model outputs.
    \item Train a reward model on preferences.
    \item Optimize the LLM policy via PPO with KL constraint.
\end{enumerate}
\end{definition}

\begin{example}[ChatGPT and Claude]
Both ChatGPT (OpenAI) and Claude (Anthropic) use RLHF as a key
training stage. The process transforms a capable but uncontrolled base
model into an assistant that follows instructions, refuses harmful
requests, and provides calibrated uncertainty.
\end{example}

\begin{definition}[DPO for Alignment]
\index{DPO}
Direct Preference Optimization eliminates the separate reward model
and RL loop, directly optimizing the policy from preference pairs. DPO
has become popular for its simplicity and stability:
\[
    \mathcal{L}_{\text{DPO}} = -\E\left[\log \sigma\left(
    \beta \log \frac{\pi_\theta(y_w|x)}{\pi_{\text{ref}}(y_w|x)}
    - \beta \log \frac{\pi_\theta(y_l|x)}{\pi_{\text{ref}}(y_l|x)}
    \right)\right]
\]
\end{definition}

\begin{remark}
Recent alternatives to RLHF include:
\begin{itemize}
    \item \textbf{DPO}: direct preference optimization (no RL loop).
    \item \textbf{RLAIF}: RL from AI feedback (using a model as judge).
    \item \textbf{Constitutional AI}: self-improvement via critique and
    revision guided by principles.
    \item \textbf{KTO}: Kahneman-Tversky Optimization (uses only
    binary good/bad labels, no pairwise comparisons).
\end{itemize}
\end{remark}

% ------------------------------------------------------------
\section{Other Applications}
\label{sec:app:other}
% ------------------------------------------------------------

\begin{example}[Science and Engineering]
\begin{itemize}
    \item \textbf{Protein folding}: AlphaFold uses RL-like techniques
    for structure prediction.
    \item \textbf{Chip design}: Google's RL-based chip placement
    achieves superhuman floorplanning.
    \item \textbf{Nuclear fusion}: DeepMind's RL controller for plasma
    confinement in tokamaks.
    \item \textbf{Compiler optimization}: MLIR/LLVM use RL for register
    allocation and instruction scheduling.
\end{itemize}
\end{example}

\begin{example}[Healthcare]
RL has been applied to:
\begin{itemize}
    \item Treatment planning (e.g., sepsis management, ventilator control)
    \item Drug discovery and molecular optimization
    \item Adaptive clinical trials
\end{itemize}
These applications require offline RL (no live experimentation on
patients) and rigorous safety constraints.
\end{example}

% ------------------------------------------------------------
\section{Challenges and Open Problems}
\label{sec:app:challenges}
% ------------------------------------------------------------

\begin{attention}[title=\textbf{Key Open Challenges}]
\begin{enumerate}
    \item \textbf{Sample efficiency}: real-world interactions are expensive.
    \item \textbf{Safety}: ensuring constraints during training and
    deployment.
    \item \textbf{Generalization}: policies that transfer across tasks
    and environments.
    \item \textbf{Reward specification}: designing rewards that capture
    true objectives (avoiding reward hacking).
    \item \textbf{Scalability}: training with billions of parameters
    and massive action spaces.
\end{enumerate}
\end{attention}

% ------------------------------------------------------------
\section{Exercises}
\label{sec:app:exercises}
% ------------------------------------------------------------

\begin{exercise}[Game Environment]
Choose an Atari game from Gymnasium and train a DQN agent from scratch.
Report the score after 1M frames and compare with human performance.
Identify which game features make RL easy or hard.
\end{exercise}

\begin{exercise}[Recommendation MDP]
Design an MDP for a music recommendation system. Define states, actions,
rewards, and transitions. Discuss the trade-off between exploration
(recommending new artists) and exploitation (playing known favorites).
\end{exercise}

\begin{exercise}[RLHF Simulation]
Simulate a simplified RLHF pipeline: (1) generate text with a small GPT,
(2) collect synthetic preferences using a rule-based scorer, (3) train a
reward model, (4) fine-tune with PPO. Measure alignment improvement on
a held-out test set.
\end{exercise}

\begin{exercise}[Sim-to-Real Gap]
Train a PPO agent for \texttt{CartPole-v1} with randomized physics
parameters (pole length, cart mass, gravity). Test on the default
parameters. Does domain randomization improve robustness?
\end{exercise}
